Regression with the Tsetlin Machine (RegressionTM) is interpretable but typically
requires densely sampled training data, raising the question of how well it can be
trained when data is scarce. We study teacher--student knowledge distillation for
RegressionTM, in which a frozen, better-informed teacher acts inside the student's
training loop: at each sample it adjusts the student's per-sample error before the
Type~I/Type~II feedback decision, steering the clause-level update directly. We develop a
family of such mechanisms and evaluate them under a masked-distillation protocol in which
nested structured holes carve progressively sparser training sets, paired by seed on
held-out data across datasets spanning the full range of teacher quality. We find that
these mechanisms modulate rather than correct the teacher's signal; that teacher accuracy
does not predict benefit; that a 70\%-data guided student can match or exceed a full-data
model where additional data is harmful; and that data spacing and density govern benefit.
Distillation in RegressionTM thus behaves as a predominantly global reshaping of the
learned function that a sufficiently accurate teacher supplements with local repair of the
regions whose training support was removed.
